\section{One positive integer weight realizes all bounded deg-lex lower terms}

\begin{theorem}[An explicit radix weight realizes bounded lexicographic comparisons]
\label{hir:actual-bounded-lex-radix-weight}
\lean{Hironaka.HironakaActualDegLexPositiveWeights.bounded_lex_radix_strict}
Let $a,b\in\mathbb N^n$, $|a|\leq D$, and $a<_{\rm lex}b$,
with the first coordinate having highest lexicographic priority.
For $B=D+2$ the literal integer sum
$L(a)=\sum_i a_i B^{n-i}$ satisfies $L(a)<L(b)$.
The upper exponent $b$ need not be bounded.
\leanok
\end{theorem}
\begin{proof}
Choose the first differing coordinate $j$: all previous coordinates
agree and $a_j<b_j$. All weights after $j$ are at most
$B^{n-j-1}$. Their total contribution for $a$ is at most
$D B^{n-j-1}<B^{n-j}$. Split each finite sum into the common
initial part, coordinate $j$, and the remaining part. Increasing
the $j$-th coefficient by one outweighs the entire remaining part
of $a$; the remaining contribution of $b$ is nonnegative.
This proves the asserted strict inequality directly in the natural
numbers, including when the remaining index set is empty.
\end{proof}

\begin{theorem}[A single positive integer weight realizes every bounded deg-lex comparison]
\label{hir:actual-deg-lex-positive-integer-weights}
\uses{hir:actual-bounded-lex-radix-weight}
\lean{Hironaka.HironakaActualDegLexPositiveWeights.exists_positive_weights_of_bounded_degLex}
For every $n,D$ there is one $w\in\mathbb N_{>0}^n$ such that
for every $a,b\in\mathbb N^n$, $|a|\leq D$ and
$a<_{\rm deglex}b$ imply $w\cdot a<w\cdot b$.
The same weight works for all such pairs, including arbitrarily
large upper exponents.
\leanok
\end{theorem}
\begin{proof}
Use the previous literal radix weights, and choose
$M=D B^n+1$. Set $w_i=M+B^{n-i}$, which is positive.
Its literal weight on an exponent is $M|a|+L(a)$.
For $|a|\leq D$, each radix factor is at most $B^n$, so
$L(a)\leq D B^n<M$. If $|a|<|b|$, one extra unit of
the common degree contribution outweighs the entire radix
contribution of $a$, while $L(b)\geq0$. If the degrees are equal,
the actual deg-lex relation is exactly the lexicographic relation,
and the previous theorem applies. Both cases use the same explicitly
constructed positive integer weight. When $n=0$ there are no strict
comparisons, so the same construction satisfies the assertion.
\end{proof}

These weights are genuine outputs from the finite degree bound. They
are intended to produce a common positive radius for actual analytic
vertex division; this finite order theorem itself makes no analytic
convergence, prepared-family or HS-history claim.
